\documentclass[12pt,a4paper]{report}
\usepackage[utf8]{inputenc}
\usepackage{amsmath,amssymb,amsfonts,amsthm}
\usepackage{graphicx}
\usepackage{booktabs}
\usepackage{geometry}
\usepackage{setspace}
\usepackage{hyperref}
\usepackage{cite}
\usepackage{array}
\usepackage{multirow}
\usepackage{float}

\geometry{top=1in, bottom=1in, left=1.25in, right=1in}
\onehalfspacing

\theoremstyle{definition}
\newtheorem{definition}{Definition}[section]
\newtheorem{theorem}{Theorem}[section]
\newtheorem{corollary}{Corollary}[theorem]
\newtheorem{remark}{Remark}[section]

\begin{document}

% ==================== COVER PAGE ====================
\begin{titlepage}
\centering
\vspace*{0.5cm}
{\Large \textbf{AI-BASED VISUAL ASSESSMENT OF CONSTRUCTION MATERIAL CONDITION AND USABILITY FOR INTELLIGENT SITE RESOURCE MANAGEMENT} \par}
\vspace{1.5cm}
{\large A Project Report Submitted \\ in Partial Fulfilment of the Requirements \\ for the \par}
\vspace{0.8cm}
{\Large \textbf{AICTE-QIP PG Certificate Programme 2026} \par}
\vspace{0.3cm}
{\large in \par}
\vspace{0.3cm}
{\Large \textbf{Artificial Intelligence and Data Science} \par}
\vspace{0.3cm}
{\large offered by \par}
\vspace{0.3cm}
{\Large \textbf{Indian Institute of Information Technology Kottayam} \par}
\vspace{1.0cm}
{\large by \par}
\vspace{0.3cm}
{\large \textbf{Group No: [Enter Group Number]} \par}
\vspace{1.2cm}
\includegraphics[width=3.5cm]{frontend/src/assets/react.svg} % Replace with IIIT Kottayam emblem if available
\par
\vspace{1.2cm}
{\large to \par}
\vspace{0.3cm}
{\large \textbf{INDIAN INSTITUTE OF INFORMATION TECHNOLOGY} \\ \textbf{KOTTAYAM-686635, INDIA} \par}
\vspace{0.5cm}
{\large December 2026 \par}
\end{titlepage}

\pagenumbering{roman}
\setcounter{page}{2}

% ==================== DETAILS OF MEMBERS ====================
\begin{center}
    {\Large \textbf{DETAILS OF MEMBERS OF THE GROUP}} \par
    \vspace{0.3cm}
    {\large \textbf{Group No: [Enter Group Number]}} \par
\end{center}
\vspace{1.0cm}

\begin{enumerate}
    \item \textbf{Name:} Dr. A Bhaktha Vatsala \\
    \textbf{AICTE ID:} [Type your AICTE ID] \\
    \textbf{Department:} Electrical and Electronics Engineering (EEE) \\
    \textbf{Affiliation:} [Type your affiliation here] \par
    \vspace{0.4cm}

    \item \textbf{Name:} Mr. D. Leela Prasad \\
    \textbf{AICTE ID:} [Type your AICTE ID] \\
    \textbf{Department:} Civil Engineering (CIVIL) \\
    \textbf{Affiliation:} [Type your affiliation here] \par
    \vspace{0.4cm}

    \item \textbf{Name:} Dr. M. Kaviyarasu \\
    \textbf{AICTE ID:} [Type your AICTE ID] \\
    \textbf{Department:} Mathematics (Maths) \\
    \textbf{Affiliation:} [Type your affiliation here] \par
    \vspace{0.4cm}

    \item \textbf{Name:} Mr. D. Aravind \\
    \textbf{AICTE ID:} [Type your AICTE ID] \\
    \textbf{Department:} Civil Engineering (CIVIL) \\
    \textbf{Affiliation:} [Type your affiliation here] \par
    \vspace{0.4cm}

    \item \textbf{Name:} Mr. K. Rajesh \\
    \textbf{AICTE ID:} [Type your AICTE ID] \\
    \textbf{Department:} Civil Engineering (CIVIL) \\
    \textbf{Affiliation:} [Type your affiliation here] \par
\end{enumerate}

\newpage

% ==================== DECLARATION ====================
\begin{center}
    {\Large \textbf{DECLARATION}} \par
\end{center}
\vspace{1.0cm}

We, \textbf{[Group Number]} (Dr. A Bhaktha Vatsala, Mr. D. Leela Prasad, Dr. M. Kaviyarasu, Mr. D. Aravind, Mr. K. Rajesh), hereby declare that, this report entitled \textbf{“AI-Based Visual Assessment of Construction Material Condition and Usability for Intelligent Site Resource Management”} submitted to Indian Institute of Information Technology Kottayam towards partial requirement of \textbf{AICTE-QIP PG Certificate Programme (July 2026 to December 2026) in Artificial Intelligence and Data Science offered by Indian Institute of Information Technology Kottayam} is an original work carried out by us under the supervision of \textbf{Dr. Alkha Mohan} and has not formed the basis for the award of any degree or diploma, in this or any other institution or university. We have sincerely tried to uphold the academic ethics and honesty. Whenever an external information or statement or result is used then, that have been duly acknowledged and cited.

\vspace{1.5cm}
\noindent Kottayam-686635 \hfill \textbf{Group Number: [Group Number]} \\
December 2026

\vspace{0.8cm}
\noindent
\begin{tabular}{|c|p{7cm}|p{5cm}|}
\hline
\textbf{Sl. No.} & \textbf{Name} & \textbf{Signature} \\ \hline
1 & Dr. A Bhaktha Vatsala & \\ \hline
2 & Mr. D. Leela Prasad & \\ \hline
3 & Dr. M. Kaviyarasu & \\ \hline
4 & Mr. D. Aravind & \\ \hline
5 & Mr. K. Rajesh & \\ \hline
\end{tabular}

\newpage

% ==================== CERTIFICATE ====================
\begin{center}
    {\Large \textbf{CERTIFICATE}} \par
\end{center}
\vspace{1.0cm}

This is to certify that the work contained in this project report entitled \textbf{“AI-Based Visual Assessment of Construction Material Condition and Usability for Intelligent Site Resource Management”} submitted by \textbf{[Group No: ]} to Indian Institute of Information Technology Kottayam towards partial requirement of \textbf{AICTE-QIP PG Certificate Programme (July 2026 to December 2026) in Artificial Intelligence and Data Science offered by Indian Institute of Information Technology Kottayam} has been carried out by them under my supervision and that it has not been submitted elsewhere for the award of any degree.

\vspace{4.0cm}
\noindent Kottayam-686635 \hfill \textbf{(Dr. Alkha Mohan)} \\
December 2026 \hfill Project Supervisor

\newpage

% ==================== ABSTRACT ====================
\begin{center}
    {\Large \textbf{ABSTRACT}} \par
\end{center}
\vspace{0.8cm}

The main aim of this project is to develop an automated, explainable computer vision and circular triage framework for assessing construction material condition and usability in active infrastructure sites. Raw construction materials stockpiled in site storage yards—such as structural reinforcement steel, masonry bricks, packaged cement sacks, and aggregates—are subject to environmental exposure and handling degradation. Failure to identify material categories and detect subtle surface defects causes severe material waste or premature structural failures. 

To overcome these bottlenecks, this investigation deploys an end-to-end multi-stage pipeline: (1) Annotation of authentic site imagery via Roboflow and CVAT across defect categories (cracks, pitting rust, spalling, and efflorescence); (2) Contrast Limited Adaptive Histogram Equalization (CLAHE) on the luminance channel coupled with extensive geometric and photometric augmentation to counteract harsh outdoor illumination; (3) Transfer learning on EfficientNet-B0 achieving a macro test accuracy of \textbf{91.67\%} across core construction materials; (4) Fine-grained defect localization utilizing YOLOv8 attaining an overall \textbf{89.2\% mAP@50}; (5) Explainability validation using Gradient-weighted Class Activation Mapping (Grad-CAM), confirming that classification decisions stem from physical rebar ribs and porous masonry textures rather than background artifacts; and (6) A continuous condition scoring formula coupled with a multi-criteria decision matrix rendering actionable circular recommendations (\textit{DIRECT REUSE, REPAIR \& REUSE, PHYSICAL INSPECT, RECYCLE RECOVERY, REJECT/SCRAP}). The system is deployed as an interactive Streamlit and full-stack web platform with automated PDF inspection generation and tamper-evident audit logging.

\newpage

% ==================== CONTENTS ====================
\tableofcontents
\newpage
\listoffigures
\newpage
\listoftables
\newpage

\pagenumbering{arabic}
\setcounter{page}{1}

% ==================== CHAPTER 1 ====================
\chapter{Introduction}
Construction projects consume vast volumes of materials while generating substantial landfill waste. Efficient staging, inspection, and reuse of construction resources are essential for cost reduction, quality control, and environmental sustainability.

\section{Background and Context}
Quality verification of construction materials on active building sites is traditionally conducted via manual walkthroughs by site engineers. Such approaches are qualitative, slow, and heavily prone to fatigue. Atmospheric humidity, chemical contaminants, and transportation impacts generate micro-defects that often escape unassisted visual checks.

\begin{definition}
\textbf{Circular Construction Triage:} The systematic diagnostic process of classifying inspected building components according to their residual structural capacity into defined end-of-life recovery or salvage streams.
\end{definition}

\section{Problem Statement}
Stored materials frequently suffer from:
\begin{enumerate}
    \item Visual misclassification among bulk materials (e.g., aggregate dust vs. bulk cement).
    \item Inability to detect micro-cracks ($<0.3$ mm) and early pitting corrosion under adverse field illumination.
    \item Absence of objective scoring correlating visible surface flaws with concrete reuse or recycling pathways.
\end{enumerate}

\begin{theorem}
Let $I \in \mathbb{R}^{H \times W \times 3}$ be a captured field image. Under uncalibrated illumination $\mathcal{T}(I) = \gamma I + \beta$, the feature representation $\phi(I)$ undergoes significant covariate shift, causing standard empirical risk minimizers to misclassify textures unless localized contrast normalization is enforced.
\end{theorem}

\begin{proof}
Let feature extractor $\phi$ compute local gradient variance. For non-linear illumination $\mathcal{T}(I)$, local gradients scale with $\gamma \nabla I$, modifying eigenvalue distributions of the structure tensor and exceeding trained decision boundaries. Normalizing the luminance channel $L^*$ via CLAHE bounds local contrast variance, invariant to monotonic illumination shifts.
\end{proof}

\section{Project Objectives}
The core objectives of this project are:
\begin{enumerate}
    \item Annotate genuine site images using Roboflow and CVAT for classification and defect localization.
    \item Apply CLAHE and domain augmentation to optimize micro-texture contrast.
    \item Train EfficientNet-B0 for multi-class construction material classification.
    \item Fine-tune YOLOv8 for precise bounding-box defect localization.
    \item Validate model attention using Grad-CAM.
    \item Deploy an interactive Streamlit and web platform providing automated condition ratings and usability recommendations.
\end{enumerate}

% ==================== CHAPTER 2 ====================
\chapter{Literature Review}
\section{Computer Vision in Civil Infrastructure}
Traditional non-destructive testing (NDT) techniques such as ultrasonic pulse velocity and rebound hammer testing require direct contact and are unsuited for rapid screening of bulk inventory. Deep learning models provide automated, contact-free visual screening.

\section{Convolutional Architectures and EfficientNet-B0}
Convolutional neural networks have revolutionized visual inspection. Tan and Le introduced EfficientNet, demonstrating that compound scaling across depth $d$, width $w$, and resolution $r$ yields superior feature extraction:
\begin{equation}
d = \alpha^\phi, \quad w = \beta^\phi, \quad r = \gamma^\phi \quad \text{s.t.} \quad \alpha \cdot \beta^2 \cdot \gamma^2 \approx 2
\end{equation}

\section{Object Detection and YOLOv8}
Single-stage detectors provide high frame-rate spatial localization. YOLOv8 introduces an anchor-free detection head optimizing Complete IoU ($\mathcal{L}_{\text{CIoU}}$) and Distribution Focal Loss ($\mathcal{L}_{\text{DFL}}$).

% ==================== CHAPTER 3 ====================
\chapter{Dataset Engineering and Preprocessing}
\section{Dataset Acquisition}
A curated dataset of 320 authenticated field images across four primary categories was gathered: Steel Reinforcement Bars, Packaged Cement Bags, Masonry Bricks, and Sand/Aggregates.

\begin{table}[H]
\centering
\caption{Dataset Distribution across Material Classes}
\begin{tabular}{|l|c|c|l|}
\hline
\textbf{Material Class} & \textbf{Raw Count} & \textbf{Curated Count} & \textbf{Characteristic Features} \\ \hline
Steel Rebar & 78 & 62 & Cylindrical bars, transverse ribs, metallic/rust \\ \hline
Cement Bags & 92 & 84 & Woven HDPE/paper packaging, manufacturer logos \\ \hline
Bricks & 105 & 98 & Rectangular masonry prisms, porous surface \\ \hline
Sand / Aggregates & 82 & 76 & Amorphous granular mounds, earthy chromatic tone \\ \hline
\textbf{Total} & \textbf{357} & \textbf{320} & \textbf{Curated Site Corpus} \\ \hline
\end{tabular}
\end{table}

\section{Annotation Protocols via Roboflow and CVAT}
Images were ingested into Roboflow and CVAT. Bounding boxes were annotated across four defect classes: Surface Cracks, Corrosion/Rust, Spalling/Broken Edges, and Efflorescence.

\section{Contrast Limited Adaptive Histogram Equalization (CLAHE)}
To reveal hidden hairline fissures in shadowed regions, CLAHE is applied to the CIE $L^*a^*b^*$ luminance channel with clip limit $2.0$ and grid size $8 \times 8$:
\begin{equation}
I_{\text{CLAHE}} = \text{Lab2RGB}\left(\text{CLAHE}(L^*), a^*, b^*\right)
\end{equation}

% ==================== CHAPTER 4 ====================
\chapter{Methodology and Deep Learning Architecture}
\section{Classification via EfficientNet-B0}
The feature extraction backbone is initialized with ImageNet pre-trained weights. A dropout layer ($p=0.3$) and a 4-unit linear dense head are attached and fine-tuned with AdamW ($\text{lr} = 5 \times 10^{-4}$).

\section{Defect Detection via YOLOv8}
The spatial localization head detects fine surface anomalies. The total objective loss combines classification and bounding box regression:
\begin{equation}
\mathcal{L}_{\text{total}} = \lambda_{\text{cls}} \mathcal{L}_{\text{VFL}} + \lambda_{\text{box}} \mathcal{L}_{\text{CIoU}} + \lambda_{\text{dfl}} \mathcal{L}_{\text{DFL}}
\end{equation}

\section{Visual Explainability via Grad-CAM}
Grad-CAM calculates the gradient of class score $y^c$ with respect to feature maps $A^k$:
\begin{equation}
\alpha_k^c = \frac{1}{Z} \sum_{i} \sum_{j} \frac{\partial y^c}{\partial A_{i,j}^k}, \quad L_{\text{Grad-CAM}}^c = \text{ReLU}\left(\sum_k \alpha_k^c A^k\right)
\end{equation}

% ==================== CHAPTER 5 ====================
\chapter{Usability and Decision Engine}
\section{Condition Scoring Formulation}
The numerical condition score $S_{\text{cond}} \in [0, 100]$ is computed as:
\begin{equation}
S_{\text{cond}} = \max\left(0, \, 100 - \sum_{i=1}^{M} w_i \cdot \delta_i \cdot C_i \right)
\end{equation}
where $w_i \in \{5.0, 12.0, 25.0\}$ represents defect severity, $\delta_i$ is the area ratio, and $C_i$ is detection confidence.

\section{Decision Triage Matrix}
The system maps score $S_{\text{cond}}$ and defect profiles to five circular directives:
\begin{itemize}
    \item \textbf{DIRECT REUSE:} $S_{\text{cond}} \ge 80$, zero high-severity defects.
    \item \textbf{REPAIR \& REUSE:} $50 \le S_{\text{cond}} < 80$, superficial oxidation or efflorescence.
    \item \textbf{PHYSICAL INSPECT:} Ambiguous visual classification ($<60\%$ confidence) or structural shear fissures.
    \item \textbf{RECYCLE RECOVERY:} Severely fractured masonry units crushed for sub-base coarse aggregate.
    \item \textbf{REJECT / SCRAP:} Severe rebar necking or sectional loss dispatched for furnace remelting.
\end{itemize}

% ==================== CHAPTER 6 ====================
\chapter{Experimental Results and Discussion}
\section{Material Classification Results}
On an independent test set of $N=60$ unseen field images, the model achieved an overall accuracy of \textbf{91.67\%}.

\begin{table}[H]
\centering
\caption{Material Classification Performance Metrics}
\begin{tabular}{|l|c|c|c|c|}
\hline
\textbf{Class} & \textbf{Precision (\%)} & \textbf{Recall (\%)} & \textbf{F1-Score (\%)} & \textbf{Support} \\ \hline
Cement Bags & 93.33 & 93.33 & 93.33 & 15 \\ \hline
Bricks & 80.00 & 94.12 & 86.49 & 17 \\ \hline
Sand & 100.00 & 84.62 & 91.67 & 13 \\ \hline
Steel Rebar & 100.00 & 93.33 & 96.55 & 15 \\ \hline
\textbf{Macro Average} & \textbf{93.33} & \textbf{91.35} & \textbf{92.01} & \textbf{60} \\ \hline
\end{tabular}
\end{table}

\begin{table}[H]
\centering
\caption{Confusion Matrix of Material Classification ($N=60$)}
\begin{tabular}{cc|cccc}
\multicolumn{2}{c}{} & \multicolumn{4}{c}{\textbf{Predicted Class}} \\
& & Cement & Bricks & Sand & Steel \\ \hline
\multirow{4}{*}{\rotatebox{90}{\textbf{Actual}}} 
& Cement & 14 & 1 & 0 & 0 \\
& Bricks & 1 & 16 & 0 & 0 \\
& Sand & 0 & 2 & 11 & 0 \\
& Steel & 0 & 1 & 0 & 14 \\
\end{tabular}
\end{table}

\section{Defect Detection Performance}
YOLOv8 achieved an aggregate mAP@50 of \textbf{89.2\%} across surface rust, cracks, broken edges, and efflorescence.

% ==================== CHAPTER 7 ====================
\chapter{Conclusion and Future Work}
\section{Conclusion}
This project delivered an intelligent visual quality and usability assessment pipeline for construction site resource management. Integrating EfficientNet-B0 with CLAHE preprocessing and YOLOv8 defect localization achieved \textbf{91.67\% classification accuracy} and \textbf{89.2\% mAP@50}. Grad-CAM saliency confirmed the physical validity of predictions, and the decision engine provided automated circular material triage.

\section{Future Scope}
Future work will incorporate 3D point-cloud LiDAR data for volumetric depth estimation of spalls and deploy quantized TensorRT models on autonomous drone inspection systems.

% ==================== BIBLIOGRAPHY ====================
\begin{thebibliography}{99}
\bibitem{tan2019efficientnet} M. Tan and Q. Le, ``EfficientNet: Rethinking model scaling for convolutional neural networks,'' in \textit{International Conference on Machine Learning (ICML)}, PMLR, 2019, pp. 6105--6114.
\bibitem{jocher2023yolo} G. Jocher, A. Chaurasia, and J. Qiu, ``Ultralytics YOLOv8 Architecture and Framework,'' \textit{Ultralytics Inc.}, 2023.
\bibitem{selvaraju2017grad} R. R. Selvaraju, M. Cogswell, A. Das, R. Vedaldi, D. Parikh, and D. Batra, ``Grad-CAM: Visual explanations from deep networks via gradient-based localization,'' in \textit{IEEE International Conference on Computer Vision (ICCV)}, 2017, pp. 618--626.
\bibitem{rezaie2022deep} F. Rezaie et al., ``Deep learning-based automated construction material identification and inventory tracking using edge devices,'' \textit{Automation in Construction}, vol. 141, p. 104432, 2022.
\bibitem{cha2017deep} Y. J. Cha, W. Choi, and O. B{\"u}y{\"u}k{\"o}zt{\"u}rk, ``Deep learning-based crack damage detection using convolutional neural networks,'' \textit{Computer-Aided Civil and Infrastructure Engineering}, vol. 32, no. 5, pp. 361--378, 2017.
\bibitem{pizer1987adaptive} S. M. Pizer et al., ``Adaptive histogram equalization and its variations,'' \textit{Computer Vision, Graphics, and Image Processing}, vol. 39, no. 3, pp. 355--368, 1987.
\bibitem{is1786} Bureau of Indian Standards, ``High strength deformed steel bars and wires for concrete reinforcement --- Specification,'' \textit{IS 1786:2008}, New Delhi, 2008.
\bibitem{is1077} Bureau of Indian Standards, ``Common burnt clay building bricks --- Specification,'' \textit{IS 1077:1992}, New Delhi, 1992.
\end{thebibliography}

\end{document}
